% !TEX root = ../main.tex
\chapter{TINJAUAN PUSTAKA}

\section{Hasil penelitian/perancangan terdahulu}
Penelitian tentang deteksi perilaku mengemudi berbahaya memakai dua pendekatan utama, yaitu pendekatan berbasis citra atau video dan pendekatan berbasis data kendaraan. Bagian ini menguraikan lima belas penelitian terdahulu yang menjadi acuan penelitian ini, diurutkan berdasarkan tahun publikasi.

Boonmee dan Tangamchit mengembangkan sistem portabel untuk mendeteksi perilaku mengemudi buruk pada bus menggunakan akselerometer dua sumbu. Sistem tersebut mengenali empat peristiwa berkendara, yaitu \textit{take off}, \textit{brake}, \textit{turn left}, dan \textit{turn right}, dengan metode berbasis aturan (\textit{rule-based}). Setiap peristiwa kemudian diberi skor dengan aturan \textit{fuzzy logic}. Hasil pengujian yang dibandingkan dengan penilaian 11 relawan menunjukkan akurasi 88,04\% untuk \textit{take off}, 87,21\% untuk \textit{brake}, 88,84\% untuk \textit{turn right}, dan 84,36\% untuk \textit{turn left}.

Carmona dkk. mengembangkan alat analisis perilaku pengemudi berbiaya rendah yang menggabungkan data sensor kendaraan dari CAN bus, \textit{inertial measurement unit} (IMU), dan GPS. Sistem tersebut mengolah sejumlah parameter dan sinyal pada domain waktu dan frekuensi untuk mendeteksi perilaku agresif secara \textit{real-time}. Lokalisasi GPS menambahkan konteks berupa informasi peta digital dan batas kecepatan. Sistem tersebut diuji pada skenario perkotaan, antarkota, dan jalan tol.

Pramunanto dkk. mengembangkan prototipe sistem pemantau perilaku pengemudi kendaraan roda empat dengan metode Na\"{\i}ve Bayes. Sistem tersebut memakai deteksi kejadian dan ekstraksi fitur untuk mengklasifikasi perilaku pengemudi ke dalam tiga kelas, yaitu defensif, normal, dan agresif. Validasi menunjukkan akurasi 86\% untuk kejadian longitudinal dan 95,8\% untuk kejadian lateral. Penelitian ini melanjutkan arah tersebut dengan mengklasifikasi perilaku per jendela waktu memakai CNN-BiLSTM.

Chauhan dkk. mengusulkan sistem \textit{deep learning} untuk mendeteksi perilaku mengemudi ugal-ugalan di jalan raya dari rekaman video. Sistem tersebut terdiri atas dua modul, yaitu ekstraksi lintasan kendaraan dan deteksi penyok (\textit{dent}) pada badan kendaraan. Lintasan kendaraan diperoleh dari koordinat objek hasil deteksi YOLOv3. Sudut antarvektor pergerakan pada tiga \textit{frame} berurutan dipakai untuk memperkirakan probabilitas perilaku ugal-ugalan. Hasil kedua modul digabungkan, dan sistem divalidasi dengan video dari YouTube.

Heo dkk. merancang sistem tertanam untuk mendeteksi perilaku mengemudi ugal-ugalan secara otonom dan \textit{real-time} menggunakan \textit{embedded GPU}, kamera, dan model \textit{deep learning} ringan. Sistem tersebut membentuk lintasan gerak kendaraan dari citra per \textit{frame} dan memanfaatkan kisi virtual (\textit{virtual grid}) pada jalan agar pelacakan tetap akurat pada video berlaju \textit{frame} rendah, yaitu 5 \textit{fps}. Evaluasi pada lintasan kendaraan nyata dan lintasan mobil kendali jarak jauh menghasilkan \textit{recall} 100\% dan presisi 67\%. Keterbatasan sumber daya komputasi pada platform tertanam menjadi tantangan utama yang dibahas pada penelitian tersebut.

Jia dkk. mengusulkan model pengenalan perilaku mengemudi abnormal yang menggabungkan \textit{Long Short-Term Memory} dan \textit{Convolutional Neural Network} (LSTM-CNN). Perilaku abnormal yang dimaksud antara lain akselerasi cepat, pengereman darurat, dan perpindahan jalur yang cepat. Menurut penelitian tersebut, sebagian besar penelitian sebelumnya memakai metode \textit{clustering} dan \textit{shallow learning} yang sulit mengenali jenis perilaku abnormal secara akurat. Jia dkk. membangun dataset dari data mengemudi kendaraan nyata dengan mendeteksi titik akselerasi dan deselerasi ekstrem melalui analisis statistik, dan model LSTM-CNN yang dilatih dengan dataset tersebut memberikan hasil yang lebih baik. Berbeda dengan penelitian tersebut, penelitian ini memakai data telemetri dari simulator dan menempatkan lapisan konvolusi sebelum lapisan BiLSTM.

Cura dkk. mengembangkan model LSTM dan CNN untuk mengklasifikasi perilaku pengemudi bus berdasarkan deselerasi, putaran mesin saat menginjak pedal, belokan, dan perpindahan jalur. Uji deselerasi, putaran mesin, dan belokan dilakukan di lintasan uji beton, sedangkan uji pindah jalur dilakukan di jalan tol beraspal. Model hanya memakai aliran data dari CAN bus kendaraan yang diolah menjadi parameter karakteristik. LSTM dan 1D-CNN menghasilkan tingkat keberhasilan yang setara, dan CNN lebih baik dalam mengenali perilaku agresif.

Rahayu membangun sistem pemantau untuk layanan rental mobil yang membaca RPM, kecepatan, beban mesin, dan suhu mesin dari OBD-II melalui Raspberry Pi 2, serta lokasi GPS dari smartphone. Sistem tersebut menilai perilaku pengemudi dengan aturan yang ditentukan dan menghasilkan tingkat kesalahan mengemudi. Pengujian menghasilkan 173 data dengan 9 pelanggaran aturan, atau tingkat kesalahan mengemudi 5,20\%.

Khan dkk. membangun perangkat \textit{plug-and-play} untuk menganalisis perilaku mengemudi secara \textit{real-time} menggunakan prosesor berdaya komputasi rendah dan sensor \textit{inertial measurement unit} (IMU) berbiaya rendah. Penelitian tersebut tidak memakai jaringan saraf tiruan karena kompleksitas komputasi dan mahalnya prosesor yang dibutuhkan. Algoritma yang diusulkan memantau percepatan longitudinal dan lateral untuk mendeteksi pengereman mendadak, belokan tajam, dan akselerasi tiba-tiba dari kondisi berhenti. Pengujian pada \textit{test bench} menunjukkan bahwa pendekatan algoritmik tersebut memiliki latensi yang cukup rendah untuk analisis perilaku mengemudi secara \textit{real-time}.

Muhamad Fadzil dkk. mengklasifikasi perilaku pengemudi dari dataset OBD-II daring yang diperoleh dari Kaggle. Fitur yang dipilih adalah putaran mesin, kecepatan kendaraan, posisi \textit{throttle}, dan beban mesin. Label aman atau agresif dibuat dengan algoritma K-Means, lalu beberapa model \textit{machine learning} dilatih untuk memprediksi label tersebut. Decision Tree mencapai akurasi rata-rata 98,72\%, dan AdaBoost menaikkannya menjadi 99,48\%. Label pada penelitian tersebut berasal dari pengelompokan data, sedangkan penelitian ini memberi label berdasarkan skenario rekaman.

Singh dan Sharma memakai data OBD-II dan pengukuran gaya-g untuk mengklasifikasi sepuluh perilaku mengemudi, antara lain \textit{weaving}, \textit{swerving}, dan \textit{sideslipping}. Model yang dipakai adalah Random Forest, AdaBoost, dan K-Nearest Neighbors, dan Random Forest mencapai akurasi tertinggi.

Thakuri dkk. mengembangkan aplikasi penghitung status ugal-ugalan (\textit{Reckless Status}) dari rekaman video jalan raya. Kendaraan dideteksi dengan YOLO dan dilacak dengan \textit{Kalman filter} untuk memperoleh data kecepatan dan lintasan. Parameter statistik kecepatan, yaitu rata-rata, standar deviasi, varians, \textit{kurtosis}, \textit{skewness}, dan \textit{peak-to-peak}, kemudian dianalisis sebagai data deret waktu dengan \textit{Recurrent Neural Network} (RNN). Model RNN tersebut mencapai akurasi 94,12\%, presisi 1,0, \textit{recall} 91,6\%, dan \textit{F1-score} 0,8979. Ambang batas yang dipakai ditentukan secara heuristik dan hanya sesuai untuk video percobaan yang direkam.

Sutikno dkk. mengembangkan sistem deteksi pengemudi ugal-ugalan berdasarkan lintasan kendaraan yang terdiri atas tiga subsistem, yaitu deteksi mobil, pelacakan mobil dengan ByteTrack, dan klasifikasi lintasan. Untuk klasifikasi lintasan, penelitian tersebut mengusulkan YOLO11n-cls yang ditingkatkan dengan blok konvolusi dan C3k2 tambahan. Model tersebut mengklasifikasikan lintasan ke dalam tiga kelas, yaitu berkendara lurus, pindah jalur halus, dan ugal-ugalan. Model usulan mencapai akurasi 0,935 dengan waktu inferensi 0,5 ms, meningkat 4,4\% dibandingkan YOLO11n-cls. Model tersebut masih lemah dalam memprediksi kelas pindah jalur halus.

Michailidis dkk. meninjau aplikasi \textit{machine learning} berbasis data OBD-II untuk efisiensi energi, pengendalian emisi, analisis perilaku pengemudi, deteksi anomali, dan keamanan siber. Menurut tinjauan tersebut, aplikasi OBD-II konvensional memakai aturan atau ambang tetap dan hanya mengenali masalah yang sudah ditentukan, sedangkan model RNN, LSTM, dan CNN mampu belajar dari data sensor berurutan. Tinjauan tersebut merangkum beberapa penelitian perilaku pengemudi berbasis OBD-II. Lee dkk. melatih LSTM, GRU, BiLSTM, dan BiGRU pada data OBD-II dari enam pengemudi di lintasan terkendali untuk empat perilaku, yaitu aman, akselerasi agresif, belokan tajam, dan pengereman keras, dan GRU mencapai akurasi 97,5\%. Song dkk. mengenali gaya mengemudi agresif, normal, dan konservatif dengan GMM semi-supervised, dan jendela 30 detik memberi hasil terbaik dengan akurasi 92\%. Tinjauan tersebut juga mencatat tantangan penerapan, yaitu dataset yang terbatas dan kurang beragam, data OBD-II yang berderau dengan laju data rendah, model \textit{deep learning} yang sering terlalu berat untuk perangkat kendaraan, serta PID khusus merek yang menyulitkan pemindahan model. Penelitian ini menjawab tantangan tersebut dengan data latih dari simulator, model berukuran kecil, dan fitur yang hanya memakai PID standar.

Sa\u{g}ba\c{s} mengembangkan kerangka kerja berbiaya rendah yang merekam seluruh data melalui aplikasi smartphone yang terhubung ke adaptor OBD-II Bluetooth. Data yang digabungkan meliputi data unit kendali mesin (ECU), sensor gerak ponsel, metadata jalan, dan sinyal audio. Perilaku mengemudi dikelompokkan ke dalam tiga kelas, yaitu tenang, normal, dan agresif, dengan jendela waktu 3, 5, 7, dan 9 detik. CNN mencapai kinerja tertinggi, dan penggabungan beberapa sumber data menaikkan akurasi klasifikasi.

Dari penelitian-penelitian tersebut, penelitian berbasis data kendaraan sebagian besar masih menilai perilaku dengan aturan atau ambang tetap, atau memberi label melalui pengelompokan data. Penelitian ini melatih model CNN-BiLSTM dengan PReLU memakai data simulator yang labelnya ditentukan oleh skenario rekaman, lalu menguji model tersebut pada data OBD-II dari mobil nyata.

\section{Teori/Konsep Dasar}
\subsection{Perilaku Mengemudi dan Keselamatan Lalu Lintas}
\subsubsection{Kecelakaan Lalu Lintas dan Faktor Penyebabnya}
Menurut WHO, kecelakaan lalu lintas masih menjadi masalah keselamatan global yang merenggut 1,16 juta jiwa setiap tahun. Mubalus membagi faktor penyebab kecelakaan lalu lintas ke dalam tiga kelompok, yaitu faktor manusia, faktor jalan, dan faktor kendaraan. Dari ketiganya, kesalahan manusia menjadi penyebab yang paling sering terjadi. Mubalus mengutip bahwa hampir 86\% kecelakaan lalu lintas disebabkan oleh kelalaian manusia, dengan faktor tersering berupa lengah, mengantuk, tidak menjaga jarak, dan melaju dengan kecepatan tinggi. Dari sisi hukum, mengemudikan kendaraan secara tidak wajar dan tanpa konsentrasi diatur dalam Pasal~283 UU No.~22 Tahun 2009. Kecepatan juga merupakan faktor risiko yang terukur. WHO menyatakan bahwa setiap kenaikan 1\% kecepatan rata-rata berkaitan dengan kenaikan 4\% risiko kecelakaan fatal.

\subsubsection{Perilaku Mengemudi Agresif/Ugal-ugalan}
\label{subsec:agresif}
Handayani dkk., yang mengutip Tasca (2000), mendefinisikan mengemudi agresif (\textit{aggressive driving}) sebagai perilaku sengaja yang cenderung meningkatkan risiko kecelakaan dan didorong oleh ketidaksabaran, kekesalan, permusuhan, atau upaya menghemat waktu. Contoh perilaku tersebut adalah berkendara secara zig-zag dan membuntuti kendaraan lain dengan jarak terlalu dekat. Thakuri dkk. mengakui bahwa istilah \textit{reckless} bersifat abstrak dan sulit didefinisikan, sehingga mereka mendefinisikannya melalui kecepatan dan parameter statistik dengan ambang batas heuristik. Penelitian ini memakai istilah ugal-ugalan untuk perilaku mengemudi agresif menurut definisi Tasca. Kedua istilah tersebut merujuk ke kelas yang sama. Definisi Tasca memuat unsur niat dan emosi yang tidak dapat diukur dari data kendaraan. Oleh karena itu, penelitian ini memakai definisi operasional: kelas ugal-ugalan adalah sesi mengemudi dengan akselerasi dan pengereman mendadak, penahanan gigi rendah, belokan tajam, dan perpindahan jalur berulang. Model mengenali perilaku agresif yang teramati dari telemetri, bukan niat pengemudi.

\subsubsection{Indikator Perilaku Ugal-ugalan}
Boonmee dan Tangamchit memakai akselerasi sumbu X untuk mengenali \textit{take off} dan \textit{brake}, serta akselerasi sumbu Y untuk mengenali gerakan ke kiri dan ke kanan. Mereka menetapkan batas kondisi tanpa aksi, yaitu akselerasi di bawah 87,5~mg pada sumbu X dan di bawah 150~mg pada sumbu Y. Satuan mg (\textit{milli-g}) berarti seperseribu percepatan gravitasi, sehingga 150~mg sama dengan sekitar 1,5~m/s$^{2}$. Mereka juga memakai \textit{connected time} untuk membedakan peristiwa kuat dari peristiwa ringan. Khan dkk. memakai akselerasi lateral untuk mengenali belokan tajam dan perpindahan jalur mendadak.

Chauhan dkk. menggolongkan kendaraan yang berpindah jalur berkali-kali dalam waktu singkat hingga membentuk lintasan menyerupai huruf S sebagai cara berkendara yang berbahaya. Perubahan kecepatan yang mendadak juga menjadi indikator perilaku ugal-ugalan, yaitu akselerasi dan deselerasi kuat yang terjadi dalam waktu berdekatan. Thakuri dkk. memakai parameter \textit{peak-to-peak} kecepatan, yaitu selisih nilai tertinggi dan terendah, untuk membedakan cara berkendara yang stabil dari cara mengemudi yang agresif. Indikator lateral seperti belokan tajam dan perpindahan jalur memerlukan sensor gerak. Singh dan Sharma menambahkan pengukuran gaya-g, dan Carmona dkk. menambahkan IMU di samping data kendaraan.

\subsection{Telemetri Kendaraan}
\subsubsection{Pengertian dan Parameter Telemetri}
Telemetri kendaraan adalah proses pengukuran dan pengiriman data parameter kendaraan secara otomatis ke penerima untuk dipantau dan dianalisis. Parameter yang umum diukur adalah kecepatan, akselerasi longitudinal dan lateral, laju rotasi kendaraan, putaran mesin, posisi gigi transmisi, serta posisi pedal gas dan pedal rem. Pada simulator, program membaca parameter-parameter tersebut langsung dari mesin fisika, sehingga tidak dibutuhkan sensor fisik. Pada mobil nyata, sebagian parameter tersebut dapat dibaca melalui port OBD-II yang dijelaskan pada Subbab~\ref{subsec:obd}.

Besaran gerak kendaraan dinyatakan pada sistem koordinat yang menempel pada badan kendaraan. Menurut Rajamani, sumbu longitudinal menunjuk ke arah depan kendaraan, sumbu lateral menunjuk ke samping, dan sumbu vertikal menunjuk ke atas. Dengan sistem koordinat ini, akselerasi longitudinal menggambarkan percepatan dan pengereman, sedangkan akselerasi lateral menggambarkan gaya samping ketika kendaraan berbelok. Kedua nilai tersebut tidak bergantung pada arah hadap kendaraan di peta.

\subsubsection{Kecepatan, Akselerasi, dan Deselerasi}
Kecepatan adalah perubahan posisi kendaraan per satuan waktu, sedangkan akselerasi adalah perubahan kecepatan per satuan waktu. Untuk dua cuplikan berurutan dengan selang waktu $\Delta t$, keduanya dihitung dengan
\begin{equation}
  v = \frac{\Delta x}{\Delta t}
  \label{eq:kecepatan}
\end{equation}
\begin{equation}
  a = \frac{\Delta v}{\Delta t}
  \label{eq:akselerasi}
\end{equation}
dengan $\Delta x$ adalah perubahan posisi dalam meter, $\Delta v$ adalah perubahan kecepatan dalam meter per detik (m/s), dan $\Delta t$ adalah selang waktu dalam detik. Satuan akselerasi adalah meter per detik kuadrat (m/s$^{2}$). Akselerasi positif berarti kendaraan bertambah cepat. Akselerasi negatif disebut deselerasi, yaitu kondisi saat kendaraan melambat, misalnya ketika pengemudi mengerem.

Pada penelitian ini, akselerasi longitudinal dihitung dari selisih kecepatan, baik pada simulator maupun pada mobil, karena adaptor OBD-II yang dipakai tidak menyediakan data akselerasi. Cara ini memperbesar derau, karena kesalahan kecil pada kecepatan dibagi dengan $\Delta t$ yang kecil. Menurut standar SAE J1979, kecepatan kendaraan pada OBD-II dikirim dalam satuan km/jam bulat. Perubahan satu satuan, yaitu 1~km/jam atau sekitar 0,28~m/s, pada $\Delta t = 0{,}2$~s sudah menghasilkan lonjakan akselerasi sekitar 1,4~m/s$^{2}$. Oleh karena itu, akselerasi hasil selisih diratakan dengan rata-rata bergerak (\textit{moving average}) atas $m$ cuplikan terakhir:
\begin{equation}
  \bar{a}_{k} = \frac{a_{k} + a_{k-1} + \dots + a_{k-m+1}}{m}
  \label{eq:moving-average}
\end{equation}
dengan $a_{k}$ adalah akselerasi hasil selisih pada cuplikan ke-$k$. Sebagai contoh, dengan $m = 3$, nilai $\bar{a}_{k}$ adalah rata-rata akselerasi dari cuplikan saat ini dan dua cuplikan sebelumnya. Kecepatan dari simulator juga dibulatkan ke km/jam bulat sebelum akselerasi dihitung, sehingga derau pada data latih menyerupai derau pada data mobil.

\subsubsection{Putaran Mesin (RPM) dan Posisi Gigi}
Putaran mesin dalam satuan RPM adalah jumlah putaran poros engkol dalam satu menit. Nilai RPM naik ketika pengemudi menekan pedal gas lebih dalam. Pada kecepatan kendaraan yang sama, RPM ditentukan oleh posisi gigi transmisi. Gigi rendah menghasilkan RPM tinggi dan tenaga dorong besar, sedangkan gigi tinggi menghasilkan RPM rendah.

Pengemudi agresif cenderung menekan pedal gas dalam sehingga mesin berputar tinggi dan transmisi tertahan pada gigi rendah. Gaya mengemudi yang halus cenderung berada pada gigi tinggi dengan RPM rendah. Hasil pendahuluan pada BAB IV menunjukkan bahwa RPM dan posisi gigi menjadi pembeda kelas terkuat. Muhamad Fadzil dkk. juga memilih putaran mesin, kecepatan, dan posisi \textit{throttle} sebagai fitur klasifikasi perilaku pengemudi dari data OBD-II. Posisi gigi hanya dapat dibaca jika mobil mendukung PID yang sesuai, sehingga penelitian ini tidak memakai posisi gigi sebagai fitur inti.

\subsubsection{Kedalaman Pedal Gas dan Pedal Rem}
Kedalaman pedal gas (\textit{throttle}) menyatakan seberapa dalam pengemudi menekan pedal akselerator, dan kedalaman pedal rem (\textit{brake}) menyatakan seberapa dalam pengemudi menekan pedal rem. Penelitian ini menyimpan kedua nilai dalam persen, dengan 0\% berarti pedal tidak ditekan dan 100\% berarti pedal ditekan penuh.

Kecepatan dan akselerasi menunjukkan hasil gerak kendaraan, sedangkan kedua pedal menunjukkan masukan dari pengemudi. Pengemudi ugal-ugalan cenderung menekan pedal gas dalam, melepasnya dengan cepat, lalu menekan pedal rem secara mendadak. Penelitian ini memakai kedalaman pedal gas sebagai fitur inti, karena besaran yang setara dapat dibaca melalui OBD-II. Kedalaman pedal rem hanya dipakai jika mobil uji menyediakannya.


\subsubsection{Telemetri pada Kendaraan Nyata melalui OBD-II}
\label{subsec:obd}
\textit{On-Board Diagnostics} II (OBD-II) adalah antarmuka diagnostik pada kendaraan yang memberi akses ke data unit kendali mesin (ECU). Permintaan data OBD-II terdiri atas byte mode dan nomor parameter (\textit{parameter identification}, PID). Mode dan PID tersebut dijelaskan dalam standar SAE J1979, dan produsen kendaraan juga dapat mendefinisikan PID sendiri. Kendaraan tidak wajib mendukung semua PID. Menurut lembar data ELM327, PID 00 pada mode 01 menunjukkan daftar PID yang didukung kendaraan dan wajib didukung oleh semua kendaraan. PID pada mode 22 umumnya milik produsen dan tidak dipublikasikan secara luas, sehingga isinya berbeda untuk setiap merek kendaraan.

ELM327 adalah penerjemah yang menjembatani port OBD-II dengan antarmuka serial standar, sehingga komputer dapat membaca data kendaraan. ELM327 dapat mendeteksi dan menerjemahkan sembilan protokol OBD secara otomatis. Penelitian ini memakai adaptor ELM327 dengan koneksi WiFi, sehingga program Python di laptop membaca data melalui jaringan nirkabel. Setiap PID diminta satu per satu, sehingga laju cuplik bergantung pada jumlah PID yang dibaca dan waktu respons kendaraan. Menurut Michailidis dkk., laju cuplik OBD-II umumnya berkisar 1 sampai 100~Hz, bergantung pada jumlah PID yang diminta, batasan protokol, dan beban ECU. Lembar data ELM327 juga mengingatkan bahwa kendaraan lama mungkin tidak dapat menerima permintaan yang terlalu cepat.

Penelitian terdahulu memakai beberapa cara untuk membaca data kendaraan nyata. Ceven dkk. membangun perangkat akuisisi dari Raspberry Pi 3 Model B+ dan modul PiCAN 2 yang membaca data ECU secara \textit{real-time} melalui konektor OBD-II dan CAN bus. Carmona dkk. menambahkan IMU dan GPS di samping data kendaraan. Sa\u{g}ba\c{s} memakai smartphone yang terhubung ke adaptor OBD-II dan sensor gerak ponsel. Malekian dkk. membaca kecepatan dan aliran massa udara dari OBD-II untuk menghitung jarak tempuh dan konsumsi bahan bakar armada kendaraan, lalu mengirimkan datanya melalui WiFi ke server.

Penelitian ini tidak memakai IMU. Fitur model dibatasi pada besaran yang tersedia di simulator dan di mobil, sesuai Tabel~\ref{tab:padanan}.

\begin{table}[H]
  \centering
  \caption{Padanan fitur inti pada simulator dan mobil}
  \label{tab:padanan}
  \small
  \begin{tabular}{|l|l|l|p{4.2cm}|}
    \hline
    \textbf{Fitur} & \textbf{Simulator} & \textbf{Mobil (mode 01)} & \textbf{Keterangan} \\
    \hline
    Kecepatan & OutGauge & PID 0D & km/jam bulat \\
    \hline
    Putaran mesin & OutGauge & PID 0C & rpm \\
    \hline
    Posisi \textit{throttle} & OutGauge & PID 11 & PID 49 (pedal akselerator) dipakai jika tersedia \\
    \hline
    Akselerasi longitudinal & selisih kecepatan & selisih kecepatan & diratakan dengan Persamaan~\ref{eq:moving-average} \\
    \hline
  \end{tabular}
\end{table}

Posisi gigi, kedalaman pedal rem, dan sudut kemudi tidak termasuk fitur inti karena ketersediaannya berbeda pada setiap mobil. Jika pemeriksaan PID menunjukkan bahwa mobil uji menyediakan salah satu besaran tersebut, dan besaran yang sama juga tersedia di simulator, besaran itu dapat ditambahkan sebagai fitur.

\subsection{Simulator BeamNG}
\subsubsection{Gambaran Umum dan Fisika Soft-body}
BeamNG adalah simulator kendaraan yang memodelkan fisika kendaraan dengan pendekatan \textit{soft-body}. Pendekatan ini memperlakukan kendaraan sebagai benda yang dapat berubah bentuk, sedangkan model \textit{rigid body} menganggap bentuk benda tetap. Pada BeamNG, kendaraan tersusun atas titik massa (\textit{node}) yang dihubungkan oleh elemen \textit{beam}. Setiap \textit{beam} bekerja seperti pegas dan peredam. Simulator menjumlahkan gaya dari semua \textit{beam} yang terhubung ke suatu \textit{node}, lalu menghitung gerak \textit{node} tersebut dengan hukum kedua Newton.

Penelitian ini memakai BeamNG karena simulator menghitung respons kendaraan terhadap pedal dan kemudi dari model fisika. Akibatnya, akselerasi, laju yaw, dan putaran mesin yang direkam mengikuti dinamika kendaraan. Simulator juga memungkinkan peneliti merekam perilaku ugal-ugalan berulang kali tanpa risiko bagi pengemudi dan pengguna jalan.

\subsubsection{BeamNG.drive, BeamNG.tech, dan Pengembangan ADAS}
\textit{Advanced Driver Assistance Systems} (ADAS) adalah sistem yang memantau kondisi kendaraan dan lingkungan untuk membantu pengemudi, misalnya peringatan tabrakan dan pengereman darurat otomatis. Pengujian ADAS di jalan nyata mahal dan berbahaya, sehingga pengembang membutuhkan lingkungan simulasi. Sistem klasifikasi perilaku mengemudi dalam penelitian ini dapat menjadi salah satu komponen ADAS, yaitu modul yang mengenali pola mengemudi berisiko dari data kendaraan.

BeamNG tersedia dalam dua versi. BeamNG.tech adalah versi untuk riset dan industri yang menyediakan antarmuka pemrograman BeamNGpy untuk mengendalikan skenario dan membaca sensor kendaraan dari Python. BeamNG.drive adalah versi umum yang memakai mesin fisika yang sama. Penelitian ini memakai BeamNG.drive karena seluruh data yang dibutuhkan dapat diperoleh melalui protokol UDP bawaan pada subbab berikutnya, tanpa lisensi riset.

\subsubsection{Akuisisi Data Telemetri Melalui Protokol UDP}
\textit{User Datagram Protocol} (UDP) adalah protokol lapisan transport yang mengirim data dalam bentuk paket tanpa membangun koneksi terlebih dahulu. UDP tidak mengirim ulang paket yang hilang dan tidak menjamin urutan paket. Sebagai gantinya, beban tambahan UDP kecil dan latensinya rendah, sehingga protokol ini cocok untuk aliran data terus-menerus seperti telemetri.

BeamNG.drive menyediakan dua protokol UDP bawaan yang diaktifkan melalui menu \textit{Options} \textgreater{} \textit{Other} \textgreater{} \textit{Protocols}:
\begin{enumerate}
  \item \textbf{OutGauge}, yaitu protokol dengan format yang sama dengan simulator Live for Speed. Satu paket berukuran 96 \textit{byte} dan memuat informasi dasbor, antara lain kecepatan (m/s), RPM, posisi gigi, kedalaman pedal gas, pedal rem, dan kopling.
  \item \textbf{MotionSim}, yaitu protokol untuk menggerakkan simulator gerak (\textit{motion platform}). Satu paket berukuran 88 \textit{byte}, diawali penanda empat karakter \texttt{"BNG1"}, lalu diikuti 21 bilangan \textit{float}: posisi, kecepatan, akselerasi tanpa gravitasi, vektor arah atas, sudut \textit{roll}--\textit{pitch}--\textit{yaw}, kecepatan sudut, dan percepatan sudut kendaraan.
\end{enumerate}
Penelitian ini memakai OutGauge sebagai sumber fitur model, karena isinya setara dengan data yang dapat dibaca melalui OBD-II. MotionSim dipakai untuk memeriksa akselerasi hasil selisih kecepatan terhadap akselerasi yang dihitung mesin fisika simulator.

Program penerima membaca kedua aliran paket dan menggabungkannya menjadi satu baris data pada setiap detak jam dengan frekuensi tetap $f_{s}$. Selang antarcuplikan dan jumlah cuplikan dari rekaman berdurasi $T$ detik adalah
\begin{equation}
  \Delta t = \frac{1}{f_{s}}, \qquad N = f_{s}\,T
  \label{eq:fs}
\end{equation}
Pada $f_{s} = 10$~Hz, selang antarcuplikan adalah 0,1~s, dan rekaman satu menit menghasilkan 600 cuplikan. Jam tetap membuat jarak antarbaris seragam walaupun waktu tiba paket UDP berubah-ubah. Karena UDP dapat kehilangan paket, jumlah cuplikan yang tersimpan dapat lebih kecil dari $N$. Program melewati baris yang sumber datanya terlambat, dan tahap pembersihan data menangani celah yang terbentuk.

\subsection{Prapemrosesan Data Deret Waktu}
\subsubsection{Pembersihan Data dan Segmentasi}
Pembersihan data adalah proses menangani data yang tidak valid agar tidak menurunkan kualitas model. Data telemetri dari UDP dapat mengandung cuplikan yang hilang, nilai yang tidak terbaca, dan nilai ekstrem akibat kejadian di luar perilaku mengemudi, misalnya tabrakan. Baris dengan nilai kosong atau tidak hingga dihapus, dan nilai ekstrem dipotong (\textit{clipping}) pada batas fisik yang masuk akal.

Penyaringan data perlu memperhatikan hubungan antara filter dan label. Jika filter yang diterapkan pada satu kelas memakai variabel yang juga menjadi masukan model, misalnya membuang baris berkecepatan tinggi hanya dari kelas normal, model dapat membedakan kelas cukup dengan mengenali batas filter tersebut. Akurasi yang diperoleh kemudian tidak mencerminkan kemampuan mengenali pola perilaku. Penyaringan yang sah berlaku sama untuk semua kelas, misalnya membuang cuplikan ketika kendaraan diam.

Cuplikan yang hilang dapat diisi dengan interpolasi linear, yaitu menarik garis lurus di antara dua cuplikan terdekat. Interpolasi hanya wajar untuk celah yang pendek. Pada celah panjang atau baris yang dibuang oleh penyaringan, nilai hasil interpolasi tidak pernah terjadi. Pendekatan yang lebih aman adalah segmentasi, yaitu memotong rekaman menjadi segmen-segmen yang barisnya berurutan, lalu mengambil setiap jendela data hanya dari dalam satu segmen. Segmen baru dimulai ketika selisih waktu antarbaris melebihi batas tertentu atau ketika ada baris yang dibuang.

\subsubsection{Penyamaan Laju Cuplik (\textit{Resampling})}
Data simulator direkam pada 10~Hz, sedangkan laju cuplik ELM327 bergantung pada jumlah PID dan waktu respons kendaraan. Model hanya mengenali pola dengan benar jika selang antarcuplikan saat pengujian sama dengan saat pelatihan. Oleh karena itu, data simulator di-\textit{resample} ke laju cuplik mobil $f_{s}$ dengan mengambil nilai pada titik waktu berjarak $1/f_{s}$ di dalam setiap segmen. Contohnya, jika $f_{s} = 5$~Hz, satu dari setiap dua cuplikan simulator dipakai. Ukuran jendela ditetapkan dalam detik, sehingga jumlah cuplikan dalam jendela tiga detik adalah $W = 3 f_{s}$.

\subsubsection{Normalisasi (Min-Max dan Z-score)}
Fitur telemetri memiliki skala yang berbeda. RPM dapat bernilai ribuan, sedangkan laju yaw hanya bernilai sepersekian radian per detik. Jika data dimasukkan tanpa normalisasi, fitur berskala besar mendominasi perhitungan gradien, sehingga pelatihan menjadi lambat dan tidak stabil.

Normalisasi \textit{min-max} mengubah setiap nilai ke rentang 0 sampai 1:
\begin{equation}
  x' = \frac{x - x_{\min}}{x_{\max} - x_{\min}}
  \label{eq:minmax}
\end{equation}
Normalisasi \textit{Z-score} mengubah nilai agar fitur memiliki rata-rata 0 dan simpangan baku 1:
\begin{equation}
  x' = \frac{x - \mu}{\sigma}
  \label{eq:zscore}
\end{equation}
dengan $x_{\min}$ dan $x_{\max}$ adalah nilai terkecil dan terbesar fitur, $\mu$ adalah rata-rata fitur, dan $\sigma$ adalah simpangan baku fitur. Sebagai contoh, RPM 4.000 dengan rata-rata 3.000 dan simpangan baku 1.000 menjadi 1, sedangkan RPM 2.000 menjadi $-1$. Nilai di bawah rata-rata selalu menjadi negatif setelah normalisasi \textit{Z-score}.

Normalisasi \textit{min-max} sensitif terhadap satu nilai ekstrem, karena nilai tersebut menentukan $x_{\max}$ atau $x_{\min}$. Oleh karena itu, penelitian ini memakai \textit{Z-score}. Nilai $\mu$ dan $\sigma$ dihitung dari data latih saja, lalu dipakai untuk data validasi, data uji, dan inferensi \textit{real-time}, sehingga informasi data uji tidak masuk ke proses pelatihan.

\subsubsection{Sliding Window}
Satu cuplikan telemetri hanya merekam kondisi kendaraan pada satu titik waktu, sehingga belum cukup untuk menilai gaya mengemudi. Pengereman mendadak atau akselerasi agresif baru terlihat dari perubahan nilai selama beberapa detik. \textit{Sliding window} memotong rekaman telemetri menjadi potongan pendek berukuran tetap yang disebut jendela. Setiap jendela memuat $W$ cuplikan berurutan dari seluruh fitur. Setelah satu jendela terbentuk, jendela bergeser sejauh $S$ cuplikan untuk membentuk jendela berikutnya. Dari satu segmen berisi $N$ cuplikan diperoleh jendela sebanyak
\begin{equation}
  N_{w} = \left\lfloor \frac{N - W}{S} \right\rfloor + 1
  \label{eq:jumlah-window}
\end{equation}
Tanda $\lfloor \cdot \rfloor$ berarti hasil pembagian dibulatkan ke bawah. Sebagai contoh, segmen berisi $N = 200$ cuplikan dengan $W = 30$ dan $S = 5$ menghasilkan $\lfloor 170/5 \rfloor + 1 = 35$ jendela. Jika $S < W$, dua jendela yang berurutan saling tumpang tindih dan berbagi $W - S$ cuplikan.

Satu jendela menjadi satu sampel masukan bagi model dan memperoleh satu label, normal atau ugal-ugalan, sesuai perilaku pengemudi saat data direkam. Ukuran jendela menentukan panjang pola berkendara yang dilihat model dalam satu kali prediksi. Nilai $W$ dan $S$ penelitian ini ditetapkan pada BAB III.

\subsubsection{Pembagian Data dan Kebocoran Data}
Penelitian ini membagi dataset menjadi data latih, data validasi, dan data uji. Data latih dipakai untuk memperbarui bobot model. Data validasi dipakai untuk memantau pelatihan, memilih \textit{epoch} terbaik, dan menyetel \textit{hyperparameter}. Data uji dipakai satu kali pada akhir untuk mengukur kinerja model.

\textit{Stratified split} membagi dataset dengan menjaga perbandingan jumlah sampel antarkelas pada ketiga bagian. Pada data deret waktu, pembagian tidak boleh dilakukan secara acak per jendela. \textit{Sliding window} menghasilkan jendela yang saling tumpang tindih, sehingga dua jendela bertetangga hampir identik. Jika keduanya jatuh ke bagian yang berbeda, model sudah melihat isi data uji saat berlatih. Kondisi ini disebut kebocoran data (\textit{data leakage}) dan membuat hasil pengujian lebih tinggi dari kemampuan sebenarnya.

Kebocoran data dicegah dengan membagi data dalam satuan blok. Setiap segmen dipotong menjadi blok yang terdiri atas banyak cuplikan berurutan. Blok-blok inilah yang dibagi ke data latih, validasi, dan uji secara terpisah per kelas, lalu jendela dibuat di dalam setiap blok. Dengan cara ini, tidak ada jendela yang berbagi cuplikan dengan jendela di bagian lain. Pengujian yang paling ketat memakai sesi mengemudi utuh yang tidak dipakai saat pelatihan, karena blok-blok dari satu sesi tetap memiliki kondisi yang mirip, seperti rute, kendaraan, dan suasana hati pengemudi.

\subsection{Deep Learning dan Convolutional Neural Network}
\subsubsection{Jaringan Saraf Tiruan}
Jaringan saraf tiruan adalah model komputasi yang terdiri atas lapisan masukan, satu atau lebih lapisan tersembunyi, dan lapisan keluaran. Setiap neuron mengalikan setiap masukan dengan bobotnya, menjumlahkan hasilnya, menambahkan bias, lalu meneruskan hasilnya melalui fungsi aktivasi $\varphi$:
\begin{equation}
  z = w_{1}x_{1} + w_{2}x_{2} + \dots + w_{n}x_{n} + b, \qquad a = \varphi(z)
  \label{eq:neuron}
\end{equation}
dengan $x_{i}$ adalah masukan, $w_{i}$ adalah bobot, $b$ adalah bias, dan $a$ adalah keluaran neuron. Pelatihan mencari nilai $w_{i}$ dan $b$ yang membuat fungsi \textit{loss} $L$, yaitu ukuran kesalahan model, sekecil mungkin. Algoritma \textit{backpropagation} menghitung seberapa besar pengaruh setiap bobot terhadap \textit{loss}, yang disebut gradien $\partial L/\partial w$. Setelah itu, \textit{gradient descent} memperbarui bobot ke arah yang menurunkan \textit{loss}:
\begin{equation}
  w_{\mathrm{baru}} = w_{\mathrm{lama}} - \eta\,\frac{\partial L}{\partial w}
  \label{eq:gd}
\end{equation}
dengan $\eta$ adalah laju pembelajaran yang menentukan besar langkah pembaruan. Pada penelitian ini, lapisan keluaran terdiri atas satu neuron dengan aktivasi sigmoid
\begin{equation}
  \mathrm{sig}(z) = \frac{1}{1 + e^{-z}}
  \label{eq:sigmoid}
\end{equation}
Sigmoid mengubah bilangan berapa pun menjadi nilai antara 0 dan 1. Keluarannya, $\hat{y}$, dibaca sebagai probabilitas bahwa suatu jendela telemetri berasal dari perilaku ugal-ugalan.

\subsubsection{Arsitektur CNN (konvolusi, pooling, fully connected)}
\textit{Convolutional Neural Network} (CNN) menyusun tiga jenis lapisan secara berurutan. Lapisan konvolusi mengekstrak fitur lokal dari masukan. Lapisan \textit{pooling} meringkas fitur tersebut. Lapisan \textit{fully connected} menggabungkan fitur dan menghasilkan keputusan klasifikasi.

Lapisan konvolusi menggeser sebuah filter kecil di sepanjang masukan. Pada setiap posisi, filter menghitung jumlah berbobot dari nilai-nilai yang tercakup, sehingga hasilnya besar ketika pola pada data mirip dengan pola pada filter. Lapisan \textit{max pooling} berukuran 2 mengambil nilai terbesar dari setiap dua langkah berurutan. Operasi ini memperpendek deret menjadi setengahnya, misalnya dari 30 langkah menjadi 15 langkah, dan tetap menyimpan respons terkuat dari setiap pola. Lapisan \textit{fully connected} bekerja seperti neuron pada Persamaan~\ref{eq:neuron}, dan lapisan terakhirnya memakai sigmoid untuk memutuskan apakah jendela termasuk normal atau ugal-ugalan.

\subsubsection{CNN 1 Dimensi untuk Data Deret Waktu}
CNN satu dimensi (1D-CNN) menggeser filter hanya sepanjang sumbu waktu. Masukannya berupa jendela berukuran $W \times F$, yaitu $W$ langkah waktu dan $F$ fitur. Untuk satu fitur dan filter berukuran 3, keluaran pada langkah ke-$t$ adalah
\begin{equation}
  y_{t} = \varphi\big(w_{1}\,x_{t-2} + w_{2}\,x_{t-1} + w_{3}\,x_{t} + b\big)
  \label{eq:konv1d}
\end{equation}
dengan $x_{t}$ adalah nilai fitur pada langkah ke-$t$, $w_{1}$ sampai $w_{3}$ adalah bobot filter, dan $b$ adalah bias. Jika masukan memiliki 11 fitur, filter memiliki tiga bobot untuk setiap fitur, dan hasil perkalian dari seluruh fitur dijumlahkan menjadi satu nilai $y_{t}$. Filter yang sama dipakai di seluruh langkah waktu, sehingga jumlah parameternya kecil dan filter dapat mengenali pola yang sama di posisi mana pun dalam jendela.

Persamaan~\ref{eq:konv1d} hanya memakai cuplikan saat ini dan dua cuplikan sebelumnya. Bentuk ini disebut konvolusi kausal (\textit{causal convolution}). Agar langkah pertama tetap memiliki keluaran, dua nilai nol ditambahkan di awal deret (\textit{padding}), sehingga panjang keluaran sama dengan panjang masukan. Sifat kausal sesuai untuk data yang diproses secara \textit{real-time}, karena model tidak memakai informasi dari masa depan.

Pada data telemetri, satu filter dapat mempelajari pola seperti kenaikan akselerasi yang cepat disertai kenaikan RPM, atau penurunan kecepatan yang tajam setelah pedal gas dilepas. Pola-pola ini mencerminkan akselerasi agresif dan pengereman mendadak, sehingga 1D-CNN sesuai untuk mengekstrak pola lokal dari jendela telemetri.

\subsubsection{Recurrent Neural Network dan Long Short-Term Memory}
Konvolusi hanya melihat pola sepanjang ukuran filternya. Perilaku mengemudi juga memiliki ketergantungan yang lebih panjang, misalnya rangkaian akselerasi keras, pengereman, lalu belokan tajam dalam beberapa detik. \textit{Recurrent Neural Network} (RNN) memproses deret langkah demi langkah dan membawa keadaan tersembunyi $\mathbf{h}_{t}$ dari langkah sebelumnya, sehingga dapat merangkum urutan kejadian. RNN biasa sulit mempelajari ketergantungan jangka panjang karena gradiennya mengecil ketika dirambatkan balik melalui banyak langkah waktu.

\textit{Long Short-Term Memory} (LSTM) yang diusulkan Hochreiter dan Schmidhuber mengatasi masalah tersebut dengan sel memori $\mathbf{c}_{t}$ yang diatur oleh tiga gerbang. Gerbang lupa $\mathbf{f}_{t}$ menentukan bagian memori lama yang dipertahankan. Gerbang masukan $\mathbf{i}_{t}$ menentukan informasi baru $\tilde{\mathbf{c}}_{t}$ yang disimpan. Gerbang keluaran $\mathbf{o}_{t}$ menentukan bagian memori yang dikeluarkan. Setiap gerbang bernilai antara 0 dan 1 karena memakai sigmoid. Memori dan keluaran LSTM diperbarui dengan
\begin{equation}
  \mathbf{c}_{t} = \mathbf{f}_{t} \odot \mathbf{c}_{t-1} + \mathbf{i}_{t} \odot \tilde{\mathbf{c}}_{t}, \qquad
  \mathbf{h}_{t} = \mathbf{o}_{t} \odot \tanh(\mathbf{c}_{t})
  \label{eq:lstm}
\end{equation}
dengan $\odot$ adalah perkalian elemen demi elemen. Persamaan pertama berarti memori baru adalah sebagian memori lama ditambah sebagian informasi baru. Karena memori diperbarui dengan penjumlahan, gradien dapat mengalir melalui banyak langkah waktu tanpa cepat menghilang. Satu lapisan LSTM dengan $u$ unit dan $d$ masukan memiliki parameter sebanyak
\begin{equation}
  N_{\mathrm{LSTM}} = 4\,\big(u\,(d + u) + u\big)
  \label{eq:param-lstm}
\end{equation}
Angka 4 berasal dari tiga gerbang dan satu calon memori $\tilde{\mathbf{c}}_{t}$, yang masing-masing memiliki bobot sendiri.

\textit{Bidirectional} LSTM (BiLSTM) yang diusulkan Schuster dan Paliwal menjalankan dua LSTM pada deret yang sama, satu dari awal ke akhir dan satu dari akhir ke awal, lalu menggabungkan keadaan tersembunyi keduanya. Dengan demikian, representasi pada setiap langkah dalam satu jendela memuat konteks sebelum dan sesudahnya. Pada klasifikasi per jendela hal ini sah, karena seluruh isi jendela sudah tersedia saat prediksi dibuat.

\subsubsection{Arsitektur Hibrida CNN-LSTM}
Arsitektur hibrida CNN-LSTM menggabungkan kelebihan kedua jenis lapisan. Lapisan konvolusi mengekstrak pola lokal pendek dari setiap beberapa langkah waktu dan memperpendek deret melalui \textit{pooling}. Lapisan LSTM kemudian merangkum urutan pola lokal tersebut menjadi satu vektor yang mewakili seluruh jendela. Pada penelitian ini, CNN mengenali kejadian sesaat seperti lonjakan pedal atau akselerasi, sedangkan BiLSTM dan LSTM mengenali urutan kejadian-kejadian itu dalam jendela tiga detik. Jia dkk. memakai gabungan LSTM dan CNN untuk mengenali perilaku mengemudi abnormal dari data kendaraan nyata. Cura dkk. menunjukkan bahwa 1D-CNN dapat mengenali perilaku agresif dari data CAN bus.

\subsubsection{Normalisasi Lapisan}
\textit{Batch normalization} dari Ioffe dan Szegedy menormalkan keluaran suatu lapisan berdasarkan rata-rata dan varians dalam satu \textit{mini-batch}, lalu menskalakannya kembali dengan dua parameter yang dipelajari. Teknik ini menstabilkan sebaran nilai yang masuk ke lapisan berikutnya sehingga pelatihan lebih cepat dan tidak mudah divergen. \textit{Layer normalization} dari Ba dkk. melakukan normalisasi serupa, tetapi menghitung rata-rata dan varians dari seluruh unit pada satu sampel. Karena tidak bergantung pada ukuran \textit{batch}, \textit{layer normalization} cocok untuk lapisan rekuren.

\subsection{Fungsi Aktivasi}
\subsubsection{ReLU dan Masalah Dying ReLU}
\textit{Rectified Linear Unit} (ReLU) meneruskan masukan positif apa adanya dan mengubah masukan negatif menjadi nol:
\begin{equation}
  f(x) = \max(0, x)
  \label{eq:relu}
\end{equation}
Sebagai contoh, $f(2) = 2$, sedangkan $f(-3) = 0$. Untuk masukan negatif, keluaran ReLU tetap nol berapa pun nilai masukannya, sehingga gradiennya juga nol. Jika masukan suatu neuron negatif untuk semua data, bobot neuron tersebut tidak lagi menerima gradien, berhenti diperbarui, dan neuron tidak ikut belajar. Kondisi ini disebut \textit{dying ReLU}.

ReLU tidak bekerja langsung pada data masukan, tetapi pada jumlah berbobot $z$ di dalam setiap neuron sesuai Persamaan~\ref{eq:neuron}. Lapisan konvolusi pertama tetap menerima nilai akselerasi negatif apa adanya, lalu bobot filter dapat mengubah tanda nilai tersebut. Risiko \textit{dying ReLU} muncul ketika $z$ suatu neuron bernilai negatif untuk hampir seluruh data latih. Data telemetri yang telah dinormalisasi dengan \textit{Z-score} tersebar di sekitar nol, sehingga banyak neuron menerima $z$ negatif pada sebagian sampel. Pada model penelitian ini, filter konvolusi yang mati berhenti mendeteksi polanya, misalnya lonjakan akselerasi, sehingga kemampuan model membedakan perilaku mengemudi berkurang.

\subsubsection{Parametric ReLU (PReLU)}
\textit{Parametric ReLU} (PReLU) yang diusulkan He dkk. mempertahankan kemiringan pada bagian negatif dan menjadikan kemiringan itu parameter yang dipelajari:
\begin{equation}
  f(x) =
  \begin{cases}
    x, & x > 0 \\
    a\,x, & x \le 0
  \end{cases}
  \label{eq:prelu}
\end{equation}
Sebagai contoh, dengan $a = 0{,}25$, masukan $-4$ menghasilkan keluaran $-1$. Koefisien $a$ diperbarui bersama bobot jaringan melalui \textit{backpropagation}, dan nilainya hanya dipengaruhi oleh masukan negatif. Model dapat memakai satu nilai $a$ untuk seluruh lapisan atau satu nilai $a$ untuk setiap kanal. Karena keluaran untuk masukan negatif tidak lagi selalu nol, neuron tetap menerima sinyal pembelajaran walaupun masukannya negatif. Penelitian ini memakai PReLU agar model menyesuaikan perlakuan terhadap nilai negatif dengan karakter data telemetri, dan agar filter tidak mati seperti pada ReLU.

Pada model penelitian ini, PReLU dipakai pada lapisan konvolusi dan lapisan \textit{fully connected} tersembunyi. Lapisan LSTM tetap memakai sigmoid dan $\tanh$ di dalam gerbangnya sesuai Persamaan~\ref{eq:lstm}. PReLU hanya menambah satu parameter per kanal, sehingga tambahan biaya komputasinya kecil dibandingkan jumlah seluruh bobot model.


\subsection{Pelatihan dan Optimisasi Model}
\subsubsection{Fungsi Loss (Binary Cross-Entropy)}
Fungsi \textit{loss} mengukur seberapa jauh prediksi model dari label sebenarnya. Selama pelatihan, model mengubah bobotnya untuk menurunkan nilai \textit{loss}. Klasifikasi perilaku normal dan ugal-ugalan merupakan masalah dua kelas, sehingga penelitian ini memakai \textit{binary cross-entropy}:
\begin{equation}
  L = -\frac{1}{N} \sum_{i=1}^{N} \Big[ y_{i}\,\log \hat{y}_{i} + (1 - y_{i})\,\log (1 - \hat{y}_{i}) \Big]
  \label{eq:bce}
\end{equation}
dengan $N$ adalah jumlah sampel, $y_{i}$ adalah label sebenarnya (1 untuk ugal-ugalan, 0 untuk normal), dan $\hat{y}_{i}$ adalah probabilitas ugal-ugalan dari sigmoid. Untuk satu sampel ugal-ugalan, hanya suku pertama yang tersisa, sehingga $L = -\log \hat{y}$. Jika model memberi $\hat{y} = 0{,}9$, nilai \textit{loss}-nya sekitar 0,11. Jika model salah dan memberi $\hat{y} = 0{,}1$, nilai \textit{loss}-nya sekitar 2,30. Kesalahan dengan keyakinan tinggi mendapat \textit{loss} terbesar, sehingga mendapat koreksi terbesar saat pembaruan bobot.

\subsubsection{Optimizer (Adam) dan Hyperparameter}
\textit{Optimizer} adalah algoritma yang memperbarui bobot jaringan berdasarkan gradien \textit{loss}. \textit{Adaptive Moment Estimation} (Adam) dari Kingma dan Ba menyimpan rata-rata gradien dari langkah-langkah sebelumnya beserta rata-rata kuadratnya. Dari kedua nilai itu, Adam menyesuaikan besar langkah pembaruan untuk setiap parameter secara terpisah. Parameter dengan gradien yang sering berganti arah memperoleh langkah yang lebih kecil, sedangkan parameter dengan gradien yang konsisten memperoleh langkah yang mendekati laju pembelajaran. Mekanisme ini menjaga pelatihan tetap stabil tanpa penyetelan laju pembelajaran yang rumit. Laju pembelajaran juga dapat diturunkan bertahap ketika \textit{loss} validasi berhenti membaik, sehingga bobot mendekati titik optimum dengan langkah yang lebih halus.

\textit{Hyperparameter} adalah nilai yang peneliti tetapkan sebelum pelatihan dan tidak dipelajari dari data. Pada penelitian ini, \textit{hyperparameter} meliputi laju pembelajaran (\textit{learning rate}), \textit{batch size}, jumlah \textit{epoch}, ukuran dan pergeseran jendela, jumlah filter dan unit LSTM, serta laju \textit{dropout}. Penelitian ini memilih nilai-nilai tersebut berdasarkan kinerja pada data validasi, sehingga data uji tetap bebas dari proses penyetelan.

\subsubsection{Overfitting, Dropout, dan Early Stopping}
\textit{Overfitting} terjadi ketika model menghafal data latih sehingga kinerjanya turun pada data baru. Tandanya, \textit{loss} latih terus turun sementara \textit{loss} validasi naik. Risiko ini besar pada penelitian dengan dataset yang direkam dari sesi mengemudi yang terbatas.

\textit{Dropout} dari Srivastava dkk. menonaktifkan sebagian neuron secara acak pada setiap iterasi pelatihan. Dengan laju \textit{dropout} $p = 0{,}3$, sekitar 30\% neuron pada lapisan tersebut dinonaktifkan pada setiap iterasi, dan keluaran neuron yang tersisa diperbesar agar total sinyalnya tetap sama. Teknik ini memaksa jaringan tidak bergantung pada satu neuron saja. \textit{Dropout} hanya aktif saat pelatihan dan tidak dipakai saat inferensi. Karena itu, \textit{loss} latih dapat lebih tinggi daripada \textit{loss} validasi tanpa menandakan kesalahan.

\textit{Early stopping} menghentikan pelatihan ketika metrik validasi tidak membaik selama $n_{p}$ \textit{epoch} berturut-turut. Nilai $n_{p}$ disebut \textit{patience}. Setelah pelatihan berhenti, bobot dari \textit{epoch} dengan metrik validasi terbaik dikembalikan sebagai bobot model.

\subsubsection{Ketidakseimbangan Kelas (Class Imbalance)}
Ketidakseimbangan kelas terjadi ketika jumlah sampel antarkelas berbeda jauh. Rasio ketidakseimbangan dihitung dengan
\begin{equation}
  \mathrm{IR} = \frac{n_{\mathrm{mayoritas}}}{n_{\mathrm{minoritas}}}
  \label{eq:ir}
\end{equation}
dengan $n_{\mathrm{mayoritas}}$ dan $n_{\mathrm{minoritas}}$ adalah jumlah sampel kelas terbanyak dan kelas tersedikit. Model yang dilatih pada data seperti ini condong memprediksi kelas mayoritas, karena kelas tersebut paling banyak menyumbang \textit{loss}. Kecondongan ke kelas ugal-ugalan membuat banyak pengemudi normal ditandai berbahaya. Kecondongan ke kelas normal membuat pengemudi ugal-ugalan lolos tanpa terdeteksi.

Pembobotan kelas mengurangi masalah ini dengan memberi bobot lebih besar pada kelas minoritas dalam \textit{loss}. Bobot kelas $c$ dihitung dengan
\begin{equation}
  w_{c} = \frac{N}{C\,n_{c}}
  \label{eq:bobot-kelas}
\end{equation}
dengan $N$ adalah jumlah seluruh sampel, $C$ adalah jumlah kelas, dan $n_{c}$ adalah jumlah sampel kelas $c$. Sebagai contoh, dataset berisi 600 jendela normal dan 400 jendela ugal-ugalan memiliki IR 1,5. Bobot kelas normal adalah $1000/(2 \times 600) \approx 0{,}83$ dan bobot kelas ugal-ugalan adalah $1000/(2 \times 400) = 1{,}25$. Metode lain adalah \textit{oversampling} kelas minoritas, \textit{undersampling} kelas mayoritas, dan augmentasi data. Karena akurasi dapat tinggi walaupun kelas minoritas gagal dikenali, penelitian ini menilai model dengan presisi, \textit{recall}, dan \textit{F1-score} per kelas, bukan hanya akurasi.

\subsection{Evaluasi Kinerja Model}
\subsubsection{Confusion Matrix}
\textit{Confusion matrix} adalah tabel yang membandingkan prediksi model dengan label sebenarnya. Penelitian ini menetapkan ugal-ugalan sebagai kelas positif dan normal sebagai kelas negatif. Tabel~\ref{tab:confusion} menunjukkan empat nilai yang dihasilkan.

\begin{table}[H]
  \centering
  \caption{\textit{Confusion matrix} untuk klasifikasi perilaku mengemudi}
  \label{tab:confusion}
  \begin{tabular}{|l|c|c|}
    \hline
     & \textbf{Prediksi ugal-ugalan} & \textbf{Prediksi normal} \\
    \hline
    \textbf{Aktual ugal-ugalan} & TP & FN \\
    \hline
    \textbf{Aktual normal} & FP & TN \\
    \hline
  \end{tabular}
\end{table}

\textit{True positive} (TP) adalah jendela ugal-ugalan yang diprediksi ugal-ugalan. \textit{True negative} (TN) adalah jendela normal yang diprediksi normal. \textit{False positive} (FP) adalah jendela normal yang diprediksi ugal-ugalan, atau alarm palsu. \textit{False negative} (FN) adalah jendela ugal-ugalan yang diprediksi normal, atau perilaku berbahaya yang tidak terdeteksi.

\subsubsection{Akurasi, Presisi, Recall, dan F1-score}
Penelitian ini menghitung empat metrik dari nilai pada \textit{confusion matrix}. Akurasi adalah proporsi prediksi yang benar dari seluruh sampel:
\begin{equation}
  \mathrm{Akurasi} = \frac{TP + TN}{TP + TN + FP + FN}
  \label{eq:akurasi}
\end{equation}
Presisi adalah proporsi prediksi ugal-ugalan yang memang ugal-ugalan:
\begin{equation}
  \mathrm{Presisi} = \frac{TP}{TP + FP}
  \label{eq:presisi}
\end{equation}
\textit{Recall} adalah proporsi perilaku ugal-ugalan yang berhasil terdeteksi:
\begin{equation}
  \mathrm{Recall} = \frac{TP}{TP + FN}
  \label{eq:recall}
\end{equation}
\textit{F1-score} adalah rata-rata harmonik presisi dan \textit{recall}:
\begin{equation}
  F_{1} = \frac{2 \times \mathrm{Presisi} \times \mathrm{Recall}}{\mathrm{Presisi} + \mathrm{Recall}}
  \label{eq:f1}
\end{equation}

Presisi rendah berarti banyak pengemudi normal ditandai ugal-ugalan. \textit{Recall} rendah berarti banyak perilaku ugal-ugalan lolos. Pada data yang tidak seimbang, akurasi saja dapat menyesatkan, sehingga penelitian ini memakai \textit{F1-score} untuk menilai keseimbangan antara presisi dan \textit{recall}.

\subsubsection{Kurva ROC dan AUC}
Keluaran model berupa probabilitas $\hat{y}$, sehingga nilai presisi dan \textit{recall} bergantung pada ambang keputusan yang dipilih. Menurut Fawcett, kurva \textit{Receiver Operating Characteristic} (ROC) menggambarkan \textit{true positive rate} terhadap \textit{false positive rate} untuk seluruh kemungkinan ambang. Luas daerah di bawah kurva tersebut (\textit{Area Under the Curve}, AUC) dapat diartikan sebagai peluang model memberi skor lebih tinggi pada satu sampel positif acak daripada pada satu sampel negatif acak. Nilai AUC 0,5 setara dengan tebakan acak, dan nilai 1 berarti kedua kelas terpisah sempurna. AUC juga dapat dihitung untuk satu fitur saja, dengan memakai nilai fitur itu sebagai skor, untuk mengukur kemampuan fitur tersebut membedakan kedua kelas tanpa bantuan fitur lain.

\subsubsection{Latensi Inferensi untuk Sistem Real-time}
Latensi inferensi adalah waktu yang dibutuhkan model sejak menerima satu jendela masukan sampai menghasilkan prediksi. Pada sistem \textit{real-time}, hasil klasifikasi harus keluar sebelum prediksi berikutnya dijadwalkan. Jika waktu pemrosesan lebih lama dari selang tersebut, prediksi tertinggal dari kondisi mengemudi yang sedang berlangsung dan antrean data terus bertambah. Latensi total mencakup waktu prapemrosesan, yaitu normalisasi dan pembentukan jendela, ditambah waktu model menghitung prediksi.

Penelitian ini mengukur latensi dengan menjalankan inferensi berulang kali lalu mengambil rata-ratanya, karena waktu satu kali eksekusi dapat dipengaruhi beban sistem operasi pada saat itu. Hasil pengukuran menunjukkan apakah model hibrida CNN-BiLSTM dengan PReLU cukup ringan untuk dipakai pada sistem pemantau perilaku mengemudi.
